% !TEX root = main.tex
\chapter{Θεωρητικό και Μαθηματικό Υπόβαθρο}
\label{chap:background}

Το κεφάλαιο αυτό παρουσιάζει το μαθηματικό υπόβαθρο στο οποίο στηρίζεται το
υπόλοιπο της εργασίας. Η Ενότητα~\ref{sec:bg-stochastic} εισάγει τις
στοχαστικές διαδικασίες και την κίνηση Brown. Η Ενότητα~\ref{sec:bg-markov}
καλύπτει τις αλυσίδες Markov και τις στάσιμες κατανομές τους. Η
Ενότητα~\ref{sec:bg-mcmc} διατυπώνει το MCMC και αναφέρει το εργοδικό θεώρημα
που διέπει όλους τους εκτιμητές που παρουσιάζονται αργότερα. Η
Ενότητα~\ref{sec:bg-langevin} εισάγει τη SDE Langevin και προσδιορίζει το
αναλλοίωτο μέτρο της. Η Ενότητα~\ref{sec:bg-gibbs} συζητά κατανομές-στόχους
Gibbs της μορφής $\pi \propto e^{-U}$ και τις συνθήκες κανονικότητας υπό τις
οποίες η διάχυση Langevin είναι καλώς ορισμένη. Η
Ενότητα~\ref{sec:bg-heavytails} δίνει τον ακριβή χαρακτηρισμό των βαριών ουρών
μέσω της αποτυχίας των εκθετικών ροπών, η οποία είναι η ιδιότητα που οδηγεί τα
περισσότερα εμπειρικά φαινόμενα που παρατηρούνται στα επόμενα κεφάλαια.

Τυπικές εγχειριδιακές παρουσιάσεις μεγάλου μέρους αυτού του υλικού είναι οι
\citet{oksendal2003stochastic} και \citet{karatzas1991brownian} για τις
στοχαστικές διαδικασίες, \citet{meyn2009markov} για τις αλυσίδες Markov,
\citet{robert2004monte} για το MCMC, και \citet{pavliotis2014stochastic} για το
πλαίσιο Langevin και Fokker--Planck.

\section{Στοχαστικές Διαδικασίες και Κίνηση Brown}
\label{sec:bg-stochastic}

Μια \emph{στοχαστική διαδικασία} σε έναν χώρο πιθανότητας $(\Omega, \mathcal{F},
\mathbb{P})$ είναι μια ακολουθία τυχαίων μεταβλητών $X = (X_t)_{t \in T}$
παραμετροποιημένη από ένα σύνολο χρόνου $T$. Οι δύο περιπτώσεις που αφορούν την παρούσα
εργασία είναι $T = \mathbb{N}$ (διακριτός χρόνος, που χρησιμοποιείται για
τις αλυσίδες Markov MCMC) και $T = [0, \infty)$ (συνεχής χρόνος, που
χρησιμοποιείται για τη SDE Langevin).

Η κανονική διαδικασία συνεχούς χρόνου είναι η \emph{κίνηση Brown}
$W = (W_t)_{t\geq 0}$ στο $\mathbb{R}^d$, που χαρακτηρίζεται από $W_0 = 0$
σχεδόν σίγουρα, ανεξάρτητες προσαυξήσεις, $W_t - W_s \sim \mathcal{N}(0, (t-s) I_d)$
για $0 \le s < t$, και συνεχείς τροχιές δείγματος
\citep{karatzas1991brownian}. Τα στοχαστικά ολοκληρώματα ως προς την κίνηση Brown
ορίζονται μέσω της κατασκευής It\^o. Μια στοχαστική διαφορική εξίσωση της μορφής
\begin{equation}
  dX_t = b(X_t)\, dt + \sigma(X_t)\, dW_t,
  \qquad X_0 = x_0,
  \label{eq:generic-sde}
\end{equation}
με μετρήσιμη ολίσθηση $b: \mathbb{R}^d \to \mathbb{R}^d$ και συντελεστή
διάχυσης $\sigma: \mathbb{R}^d \to \mathbb{R}^{d \times d}$, έχει μοναδική ισχυρή
λύση υπό τυπικές παραδοχές Lipschitz και γραμμικής αύξησης για τα $b$ και
$\sigma$ \citep{oksendal2003stochastic}.

Για το \eqref{eq:generic-sde}, η κατανομή του $X_t$ έχει πυκνότητα $p_t(\cdot)$
που εξελίσσεται σύμφωνα με την εξίσωση \emph{Fokker--Planck},
\begin{equation}
  \partial_t p_t \;=\; -\nabla \cdot (b\, p_t) \;+\; \tfrac12 \nabla^2 : (a\, p_t),
  \qquad a = \sigma\sigma^\top,
  \label{eq:fokker-planck}
\end{equation}
με αρχική συνθήκη $p_0 = \delta_{x_0}$ (ή, πιο γενικά, την αρχική κατανομή του
$X_0$). Ένα μέτρο $\mu$ στο $\mathbb{R}^d$ είναι \emph{αναλλοίωτο} για τη SDE
αν το $\mu$ αποτελεί στάσιμη λύση του \eqref{eq:fokker-planck}, δηλ.\
$\mathcal{L}^* \mu = 0$ όπου $\mathcal{L}^*$ είναι ο συζυγής του γεννήτορα
\citep{pavliotis2014stochastic}.

\section{Αλυσίδες Markov και Στάσιμες Κατανομές}
\label{sec:bg-markov}

Μια αλυσίδα Markov διακριτού χρόνου $(X_k)_{k\geq 0}$ σε έναν μετρήσιμο χώρο
καταστάσεων $(E, \mathcal{E})$ ορίζεται από μια αρχική κατανομή και έναν πυρήνα
μετάβασης $P : E \times \mathcal{E} \to [0,1]$, όπου
$P(x, A) = \mathbb{P}(X_{k+1} \in A \mid X_k = x)$. Ο πυρήνας δρα σε φραγμένες
μετρήσιμες συναρτήσεις $\varphi$ από τα δεξιά,
\begin{equation}
  P\varphi(x) \;=\; \int_E \varphi(y)\, P(x, dy),
\end{equation}
και σε μέτρα από τα αριστερά, $(\mu P)(A) = \int_E P(x, A)\, \mu(dx)$.

Ένα μέτρο πιθανότητας $\pi$ σε $(E, \mathcal{E})$ είναι \emph{στάσιμο} για τον
$P$ αν $\pi P = \pi$. Η στασιμότητα προκύπτει από την ισχυρότερη συνθήκη
\emph{λεπτομερούς ισορροπίας}
\begin{equation}
  \pi(dx)\, P(x, dy) \;=\; \pi(dy)\, P(y, dx),
  \qquad x, y \in E,
  \label{eq:detailed-balance}
\end{equation}
η οποία αποτελεί την αρχή σχεδιασμού των αλγορίθμων Metropolis--Hastings.

Μια αλυσίδα Markov λέγεται \emph{μη υποβιβάσιμη} αν από οποιαδήποτε κατάσταση
μπορεί να φτάσει σε οποιαδήποτε άλλη με θετική πιθανότητα, \emph{απεριοδική}
αν δεν εμφανίζει κυκλική συμπεριφορά, και \emph{γνησίως επαναληπτική} αν επιστρέφει με πιθανότητα $1$ σε κάθε σύνολο θετικής μέτρησης.
Μια αλυσίδα που είναι γνησίως επαναληπτική και απεριοδική λέγεται
\emph{εργοδική}. Κάθε εργοδική αλυσίδα έχει \emph{μοναδική} στάσιμη κατανομή
$\pi$ στην οποία η περιθώρια κατανομή $\mathbb{P}(X_k \in \cdot)$ συγκλίνει
καθώς $k \to \infty$, για οποιαδήποτε αρχική κατανομή
\citep[Theorem~13.0.1]{meyn2009markov}.

\section{MCMC και το Εργοδικό Θεώρημα}
\label{sec:bg-mcmc}

Στο Markov chain Monte Carlo επιλέγουμε έναν πυρήνα μετάβασης $P$ του
οποίου η μοναδική στάσιμη κατανομή είναι η κατανομή-στόχος $\pi$, τρέχουμε
την αλυσίδα για να παράγουμε μια τροχιά $X_0, X_1, \ldots, X_n$, και
χρησιμοποιούμε τον χρονικό μέσο για να προσεγγίσουμε ολοκληρώματα ως προς την $\pi$. Το
αποτέλεσμα που δικαιολογεί αυτή την προσέγγιση είναι το \emph{εργοδικό θεώρημα}:
αν η $\pi$ είναι η μοναδική στάσιμη κατανομή μιας εργοδικής αλυσίδας και
$\varphi \in L^1(\pi)$, τότε για $\pi$-σχεδόν κάθε σημείο εκκίνησης,
\begin{equation}
  \frac{1}{n} \sum_{k=1}^n \varphi(X_k) \;\xrightarrow{a.s.}\;
  \int_E \varphi(x)\, \pi(dx) \;=\; \mathbb{E}_\pi[\varphi(X)]
  \qquad \text{καθώς } n \to \infty,
  \label{eq:ergodic-thm}
\end{equation}
\citep[Theorem~17.0.1]{meyn2009markov}. Ένα κεντρικό οριακό θεώρημα ισχύει υπό
πρόσθετες συνθήκες κανονικότητας \citep{jones2004markov, geyer1992practical}.
Η ασυμπτωτική διακύμανση είναι $\sigma^2_\varphi = \mathrm{Var}_\pi(\varphi(X_0))
\cdot \tau_\varphi$ όπου ο ολοκληρωμένος χρόνος αυτοσυσχέτισης
\begin{equation}
  \tau_\varphi \;=\; 1 + 2\sum_{k=1}^\infty \mathrm{Corr}_\pi(\varphi(X_0), \varphi(X_k))
  \label{eq:iat}
\end{equation}
δείχνει πόση πληροφορία χάνεται λόγω συσχέτισης σε σχέση με ανεξάρτητα δείγματα.
Το \emph{αποτελεσματικό μέγεθος δείγματος} $N_\mathrm{eff} = n / \tau_\varphi$,
που χρησιμοποιείται σε όλη την εργασία, προκύπτει από το \eqref{eq:iat} και
εκτιμάται στην πράξη με τον εκτιμητή αρχικής θετικής ακολουθίας του Geyer
εφαρμοσμένο στην εμπειρική αυτοσυσχέτιση
\citep{geyer1992practical}.

\section{SDE Langevin και τα Αναλλοίωτα Μέτρα τους}
\label{sec:bg-langevin}

Η \emph{overdamped SDE Langevin} στο $\mathbb{R}^d$ που συσχετίζεται με ένα
δυναμικό $U \in C^1(\mathbb{R}^d; \mathbb{R})$ είναι
\begin{equation}
  dX_t \;=\; -\nabla U(X_t)\, dt \;+\; \sqrt{2}\, dW_t.
  \label{eq:overdamped-bg}
\end{equation}
Εισάγοντας $b = -\nabla U$ και $\sigma = \sqrt{2}\, I_d$ στην εξίσωση
Fokker--Planck \eqref{eq:fokker-planck} και αναζητώντας στάσιμη λύση, προκύπτει
\begin{equation}
  0 \;=\; \nabla \cdot (\nabla U \cdot p_\infty) + \Delta p_\infty,
\end{equation}
της οποίας η μοναδική λύση πυκνότητας πιθανότητας υπό ήπιες συνθήκες αύξησης
του $U$ είναι
\begin{equation}
  p_\infty(x) \;=\; Z^{-1} \exp(-U(x)),
  \qquad Z = \int_{\mathbb{R}^d} \exp(-U(y))\, dy < \infty.
  \label{eq:invariant-overdamped}
\end{equation}
Δηλαδή, η SDE \eqref{eq:overdamped-bg} \emph{έχει το μέτρο Gibbs $\pi \propto
e^{-U}$ ως αναλλοίωτη κατανομή της}, που είναι η ιδιότητα που την καθιστά
χρήσιμη για δειγματοληψία. Η γεωμετρική εργοδικότητα ως προς την $\pi$ ισχύει
όταν το $U$ είναι έντονα κυρτό στο άπειρο
\citep{roberts1996exponential, pavliotis2014stochastic}.

Η \emph{underdamped SDE Langevin} επαυξάνει τη θέση $X_t$ με μια ορμή $V_t$ που
υπακούει σε
\begin{equation}
  \begin{cases}
    dX_t = V_t\, dt, \\
    dV_t = -\nabla U(X_t)\, dt - \gamma V_t\, dt + \sqrt{2\gamma}\, dW_t,
  \end{cases}
  \label{eq:underdamped-bg}
\end{equation}
όπου $\gamma > 0$ είναι ο συντελεστής τριβής. Το μοναδικό αναλλοίωτο μέτρο στον
κοινό χώρο $(X, V) \in \mathbb{R}^{2d}$ είναι το μέτρο Boltzmann--Gibbs
\begin{equation}
  \tilde\pi(x, v) \;\propto\; \exp\!\left(-U(x) - \tfrac12 \|v\|^2\right),
  \label{eq:invariant-underdamped}
\end{equation}
του οποίου η $x$-περιθώρια κατανομή είναι το επιθυμητό $\pi$
\citep{leimkuhler2015molecular}. Η συνιστώσα ορμής λειτουργεί επομένως ως
βοηθητική μεταβλητή που απορρίπτεται μετά τη δειγματοληψία. Σε σύγκριση με την
overdamped SDE \eqref{eq:overdamped-bg}, η underdamped διατύπωση μπορεί να
αναμειγνύεται ταχύτερα σε ορισμένα καθεστώτα επειδή η αλυσίδα θυμάται σε ποια κατεύθυνση κινούνταν, μέσω της ορμής. Βλ.\ \citet{cheng2018underdamped}
για μη-ασυμπτωτική ανάλυση σε log-κοίλες κατανομές.

\section{Κατανομές-Στόχοι Gibbs}
\label{sec:bg-gibbs}

Οι κατανομές-στόχοι που εξετάζονται στην εργασία είναι όλες Gibbs
μορφής \eqref{eq:gibbs-intro}. Εξετάζουμε τρία συγκεκριμένα δυναμικά.

\paragraph{Γκαουσιανή.}
$U(x) = \tfrac12 x^\top \Sigma^{-1} x$ για θετικά ορισμένο $\Sigma$ παράγει μια
Γκαουσιανή κατανομή-στόχο $\mathcal{N}(0, \Sigma)$. Το δυναμικό είναι έντονα
κυρτό με σταθερά ίση με τη μικρότερη ιδιοτιμή του $\Sigma^{-1}$, οπότε η
overdamped SDE Langevin είναι γεωμετρικά εργοδική και οι διακριτοποιημένοι
αλγόριθμοι διαθέτουν ανεπτυγμένη μη-ασυμπτωτική θεωρία
\citep{dalalyan2017theoretical, durmus2017nonasymptotic}.

\paragraph{Student-$t$.}
Μια κατανομή Student-$t$ με $\nu$ βαθμούς ελευθερίας, θέση $0$ και κλίμακα $1$
έχει δυναμικό
\begin{equation}
  U(x) \;=\; \frac{\nu + d}{2}\, \log\!\left(1 + \frac{\|x\|^2}{\nu}\right),
\end{equation}
του οποίου η κλίση είναι $\nabla U(x) = \frac{\nu+d}{\nu}\,\frac{x}{1 + \|x\|^2/\nu}$
και φθίνει προς το μηδέν καθώς $\|x\| \to \infty$ (η επαναφορική ολίσθηση
εξασθενεί), αντί να αυξάνεται γραμμικά όπως στη Γκαουσιανή. Το δυναμικό
αυξάνεται λογαριθμικά και όχι τετραγωνικά, άρα το $U$
\emph{δεν} είναι έντονα κυρτό στο άπειρο. Ροπές τάξης $p < \nu$ υπάρχουν,
ενώ ροπές τάξης $p \geq \nu$ δεν υπάρχουν. Ειδικότερα, η διακύμανση
υπάρχει για $\nu > 2$ και ισούται με $\nu / (\nu - 2)$. Καθώς $\nu \to \infty$, η
κατανομή Student-$t$ συγκλίνει στην τυπική Γκαουσιανή. Για $\nu = 1$ ταυτίζεται
με την τυπική κατανομή Cauchy.

\paragraph{Cauchy.}
Η τυπική κατανομή Cauchy έχει δυναμικό
\begin{equation}
  U(x) \;=\; \frac{d+1}{2}\, \log\!\left(1 + \|x\|^2\right),
\end{equation}
κλίση $\nabla U(x) = (d+1) x / (1 + \|x\|^2)$, και \emph{καμία πεπερασμένη ροπή
οποιασδήποτε θετικής τάξης}. Η ολίσθηση στην overdamped SDE \eqref{eq:overdamped-bg}
εξαφανίζεται επομένως καθώς $\|x\| \to \infty$. Η SDE εξακολουθεί να είναι
εργοδική, αλλά τα κλασικά κριτήρια γεωμετρικής εργοδικότητας δεν εφαρμόζονται,
και τα τυπικά φράγματα σφάλματος διακριτοποίησης για ULA και MALA επιδεινώνονται
σημαντικά \citep{jarner2007convergence}.

\section{Κατανομές με Βαριές Ουρές}
\label{sec:bg-heavytails}

Λέμε ότι μια πραγματικής τιμής τυχαία μεταβλητή $X$ έχει \emph{βαριές ουρές}
(με την έννοια που αφορά αυτή την εργασία) αν η γεννήτρια συνάρτηση
ροπών της δεν υπάρχει σε κανένα μη-τετριμμένο διάστημα, δηλαδή
\begin{equation}
  \mathbb{E}\!\left[\exp(\lambda X)\right] \;=\; +\infty
  \qquad \text{για κάθε } \lambda > 0.
  \label{eq:heavy-tail-def}
\end{equation}
Αυτό ισοδυναμεί με το να λέμε ότι η ουρά $\mathbb{P}(|X| > t)$ φθίνει
υπο-εκθετικά \citep[Ch.~3]{resnick2007heavy}. Μια κανονική κλάση κατανομών με
βαριές ουρές είναι η κλάση \emph{κανονικά μεταβαλλόμενων} κατανομών, για τις
οποίες $\mathbb{P}(|X| > t) = t^{-\alpha} \ell(t)$ με $\alpha > 0$ και $\ell$ μια
αργά μεταβαλλόμενη συνάρτηση. Η κατανομή Cauchy ανήκει σε αυτή την κλάση με
$\alpha = 1$, και η Student-$t$ με $\nu$ βαθμούς ελευθερίας ανήκει με $\alpha = \nu$.

Αυτές οι ιδιότητες έχουν συνέπειες για τους αλγορίθμους Langevin. Η γεωμετρική
εργοδικότητα της overdamped SDE Langevin απαιτεί τυπικά μια ανισότητα
Foster--Lyapunov της μορφής $\mathcal{L}\Psi \leq -c\Psi + b\mathbf{1}_C$ για
κάποια συνάρτηση Lyapunov $\Psi$ που αυξάνεται τουλάχιστον εκθετικά ως προς
$\|x\|$ \citep{meyn2009markov, mattingly2002ergodicity}. Για κατανομές με βαριές
ουρές, καμία τέτοια $\Psi$ δεν υπάρχει στο $L^1(\pi)$, οπότε η θεωρία δίνει
μόνο υπο-γεωμετρικούς ρυθμούς σύγκλισης \citep{jarner2007convergence}. Στην
πράξη αυτό φαίνεται ως μεγάλοι χρόνοι αυτοσυσχέτισης, αργή εξερεύνηση ουράς
και --- για τον ULA --- μόνιμη μεροληψία διακριτοποίησης.

\paragraph{Ελαφρύ, μεσαίο και βαρύ καθεστώς στην εργασία.}
Υιοθετούμε την ακόλουθη άτυπη ταξινόμηση καθ' όλα τα εμπειρικά κεφάλαια:

\begin{itemize}
  \item \textbf{Ελαφριές ουρές:} $\mathbb{P}(|X|>t)$ φθίνει τουλάχιστον
        εκθετικά. Ισοδύναμα, η MGF υπάρχει σε μια γειτονιά του μηδενός.
        Η Γκαουσιανή είναι το κανονικό παράδειγμα.
  \item \textbf{Μεσαίες ουρές:} πολυωνυμική φθορά με πεπερασμένη διακύμανση,
        δηλ.\ $\mathbb{P}(|X|>t) \sim c t^{-\alpha}$ με $\alpha > 2$.
        Η Student-$t$ με $\nu = 5$ που χρησιμοποιείται στα πειράματα έχει
        $\alpha = 5$.
  \item \textbf{Βαριές ουρές:} πολυωνυμική φθορά χωρίς πεπερασμένο μέσο,
        δηλ.\ $\alpha \leq 1$. Η Cauchy είναι η οριακή περίπτωση $\alpha = 1$.
\end{itemize}

Η ταξινόμηση αυτή είναι απλοποιημένη αλλά χρήσιμη: βοηθά να καταλάβουμε ποιες
εγγυήσεις σύγκλισης ισχύουν σε κάθε περίπτωση (όλες για ελαφριές ουρές, αργές
αλλά έγκυρες για μεσαίες, με σοβαρά προβλήματα σύγκλισης για βαριές) και
οργανώνει τα αποτελέσματα του Κεφαλαίου~\ref{chap:experiments}.

\newpage
\section{Σύνοψη Συμβολισμού}
\label{sec:bg-notation}

Συγκεντρώνουμε τον συμβολισμό που χρησιμοποιείται στο υπόλοιπο της εργασίας.

\begin{itemize}
  \item $\pi$ δηλώνει την κατανομή πιθανότητας-στόχο στο $\mathbb{R}^d$.
        $U$ δηλώνει το δυναμικό της, έτσι ώστε $\pi(x) \propto \exp(-U(x))$.
  \item $\nabla U$ δηλώνει την κλίση του δυναμικού. Στην πράξη οι αλγόριθμοι
        χρησιμοποιούν αυτό το αντικείμενο για τον υπολογισμό της $\nabla U(x)$.
  \item $V_t \in \mathbb{R}^d$ δηλώνει τη μεταβλητή ορμής της underdamped SDE
        Langevin.
  \item $\eta > 0$ δηλώνει το βήμα (ή βήμα ολοκληρωτή) της διακριτοποιημένης
        SDE Langevin.
  \item $\gamma > 0$ δηλώνει τον συντελεστή τριβής της underdamped SDE.
  \item $d$ δηλώνει τη διάσταση του $\pi$.
  \item $X_k$ δηλώνει το $k$-οστό δείγμα της αλυσίδας διακριτού χρόνου.
        $X_t$ (με συνεχή δείκτη) δηλώνει την αντίστοιχη λύση SDE συνεχούς
        χρόνου όταν το πλαίσιο καθιστά τη διάκριση σαφή.
  \item $N_\mathrm{eff}$ και $\tau$ δηλώνουν το αποτελεσματικό μέγεθος
        δείγματος και τον ολοκληρωμένο χρόνο αυτοσυσχέτισης, που συνδέονται
        με $N_\mathrm{eff} = n / \tau$ για αλυσίδα μήκους $n$.
\end{itemize}
